\chapter*{Abstract}
\addcontentsline{toc}{chapter}{Abstract}
\markboth{Abstract}{Abstract}

The accelerating integration of photovoltaic (PV) generation and battery energy storage systems (BESS) into distribution networks is displacing conventional synchronous generation, reducing system inertia and exposing microgrids to increased operational volatility. Existing rule-based and single-horizon optimisation controllers struggle to reconcile short-term forecast uncertainty with long-term degradation and reliability objectives, particularly during islanded operation following upstream network disturbances.

This thesis develops and validates an adaptive model predictive control (MPC) framework for resilient microgrid energy management under deep renewable penetration. The framework couples a rolling-horizon optimisation core with an online forecast-error correction layer, allowing the controller to reweight state-of-charge (SOC) reserve margins in response to observed prediction drift rather than relying on static safety factors. A reduced-order thermal-electrochemical battery model is embedded within the optimisation to jointly minimise operating cost and cumulative degradation.

The proposed controller is evaluated against baseline rule-based and static-horizon MPC strategies using twelve months of half-hourly irradiance, load, and tariff data drawn from a synthetic feeder representative of regional Australian conditions. Across all evaluated scenarios, the adaptive controller reduces expected daily operating cost by 14.2\% relative to the rule-based baseline and by 6.8\% relative to static-horizon MPC, while reducing SOC constraint violations during islanding events by more than half. Computation time remains within the bounds required for five-minute dispatch cycles on embedded controller hardware.

These results indicate that explicitly modelling and correcting for forecast uncertainty within the control loop -- rather than treating it as an external disturbance -- yields measurable gains in both economic performance and operational resilience. The thesis contributes a generalisable adaptive-horizon formulation, an open evaluation methodology for islanding resilience, and a set of design guidelines for practitioners deploying MPC-based controllers on resource-constrained microgrid hardware.

% ============================================================
% DECLARATION
% ============================================================
